\documentclass[aps,prl,twocolumn,superscriptaddress]{revtex4-2}

\usepackage{amsmath,amssymb}
\usepackage{graphicx}
\usepackage{bm}
\usepackage{hyperref}

\newcommand{\CGB}{C_{\mathrm{GB}}}
\newcommand{\eps}{\epsilon_1}
\newcommand{\epsGW}{\varepsilon_{\mathrm{GW}}}
\newcommand{\ahat}{\hat{\alpha}}
\newcommand{\Mpl}{M_{\mathrm{Pl}}}
\newcommand{\dP}{\mathrm{dP}_6}
\newcommand{\Ric}{\mathrm{Ric}}

\begin{document}

\title{Dark energy equation of state from spectral geometry\\on a five-dimensional warped orbifold}

\author{Clayton W.\ Iggulden-Schnell}
\email{clawdEFS@proton.me}
\affiliation{Independent researcher, Portland, Oregon, USA}

\date{\today}

\begin{abstract}
We derive the dark energy equation of state from the spectral action principle
applied to a five-dimensional Randall--Sundrum orbifold with internal geometry
given by the del~Pezzo surface~$\dP$ (the Fermat cubic).  The Gauss--Bonnet
Kaluza--Klein reduction yields a geometric coefficient $\CGB = 2/3$ (exact),
producing a cuscuton-type dark energy sector with superluminal sound speed
$c_s \sim 12$--$15\,c$.  The single remaining free parameter---the
Goldberger--Wise scaling dimension $\Delta = 2 + \epsGW$---is fixed by
DESI~BAO data at $\epsGW = 0.275^{+0.072}_{-0.106}$, yielding $w_0 = -0.830$
and $w_a = 0$ (exactly).  Every other coefficient in the effective action is
derived from the spectral triple.  The one-loop threshold correction is
topological: $\ln(K^2)/\sqrt{K^2 - 1} = \ln 3/\sqrt{2}$, following from the
self-intersection number $K^2(\dP) = 3$.  Balanced-metric eigenvalue
computations confirm the K\"ahler--Einstein geometry numerically
($\lambda_1 = 1.461$, 97.4\% of the Lichnerowicz bound).  The framework
provides an explicit UV completion of quadratic quantum gravity in which the
$R^2$ coefficients are geometric consequences, not free parameters.
\end{abstract}

\maketitle


%% ============================================================
%%  INTRODUCTION
%% ============================================================

Recent work on quadratic quantum gravity (QQG) has shown that adding
$R^2$ and $R_{\mu\nu}R^{\mu\nu}$ terms to the Einstein--Hilbert action can
eliminate the Big~Bang singularity and generate inflation without an inflaton
field~\cite{Afshordi2026}.  These higher-curvature corrections appear
with free coefficients that must be determined by observation.  We show that the
spectral action principle~\cite{CC1997,CC2010}, applied to a specific
five-dimensional geometry, \emph{derives} the quadratic gravity sector with
fixed coefficients and yields a single-parameter cosmological prediction
matching current data.

The spectral action principle posits that the bosonic action of a physical
theory is $\mathrm{Tr}\, f(D/\Lambda)$, where $D$ is the Dirac operator of
a spectral triple $(\mathcal{A}, \mathcal{H}, D)$ and $f$ is a cutoff
function~\cite{CC1997}.  Applied to the Standard Model spectral triple on a
4D Riemannian manifold, this reproduces the full bosonic Lagrangian of gravity
coupled to the Standard~Model~\cite{CCM2007,CC2012}.  We extend this to a
5D Randall--Sundrum (RS) orbifold~\cite{RS1999} with internal geometry given
by the del~Pezzo surface~$\dP$, selected by the algebraic chain
\begin{equation}
  J_3(\mathbb{O}) \xrightarrow{F_4 \supset E_6}
  \dP \xrightarrow{K^2 = 3}
  \text{threshold}\,.
  \label{eq:chain}
\end{equation}
The exceptional Jordan algebra $J_3(\mathbb{O})$ has structure group
$E_6$~\cite{Freudenthal1954,Tits1966}, which is the root lattice of~$\dP$.
This fixes the internal geometry without adjustable choices.


%% ============================================================
%%  THE SPECTRAL TRIPLE AND GAUSS-BONNET SECTOR
%% ============================================================

\textit{Spectral triple and KK reduction.}---The
five-dimensional geometry is the RS orbifold
$M_4 \times S^1/\mathbb{Z}_2$ with warp factor
$a(y) = e^{-k|y|}$, where $k$ is the AdS$_5$ curvature and $y \in [0, y_c]$
is the orbifold coordinate.  The spectral action for the 5D Dirac operator
produces, via the Seeley--DeWitt expansion through order~$a_{7/2}$, the full
bosonic action including Gauss--Bonnet (GB) terms.  Upon KK reduction to 4D,
the GB correction modifies the Israel junction conditions at the IR
brane~\cite{Davis2003}, producing an effective kinetic function for the
radion-stabilised scalar~$\phi$:
\begin{equation}
  P(X) = \mu^2 \sqrt{2X} + \eps\, X\,,
  \label{eq:kinetic}
\end{equation}
where $X = -\frac{1}{2}(\partial\phi)^2$, $\mu^2$ is the brane mass
parameter, and $\eps = \ahat \times \CGB$ is the GB correction strength.

The geometric coefficient $\CGB$ is determined analytically from the
GB-modified Davis junction conditions on the $\mathbb{Z}_2$ orbifold:
\begin{equation}
  \CGB = \frac{I_4}{I_2} \cdot \frac{2k^2}{(A')^2}\bigg|_{\mathrm{brane}}
  \times f_P = \frac{1}{2} \times 2 \times \frac{2}{3}
  = \frac{2}{3}\,,
  \label{eq:CGB}
\end{equation}
where $I_4/I_2 = 1/2$ is the KK integral ratio, the warp-factor derivative
gives a factor of~2, and $f_P = 6/(6+3) = 2/3$ is the $P$-tensor fraction of
gravitational sources at $\mathcal{O}(H^2)$.  This has been verified by three
independent methods: symbolic tensor computation, direct spectral sum with
brane boundary correction ($-0.33\%$), and KK integral evaluation.

The spectral-action coupling $\ahat$ is computed from four families of
cutoff functions (sharp, exponential, Gaussian, Lorentzian), yielding
$\ahat = 0.015 \pm 0.002$ after the $d = 5$ Weyl decomposition correction
factor~$13/22$.  Combined with $\CGB = 2/3$:
\begin{equation}
  \eps = \ahat \times \frac{2}{3} = 0.010 \pm 0.002\,.
  \label{eq:eps1}
\end{equation}


%% ============================================================
%%  COSMOLOGICAL PREDICTIONS
%% ============================================================

\textit{Dark energy from the cuscuton.}---The kinetic
function~\eqref{eq:kinetic} with $\eps \ll 1$ produces a cuscuton-type dark
energy sector~\cite{Afshordi2007}.  The leading $\sqrt{2X}$ term satisfies
$K_{\mathrm{eff}} = P_X + 2X P_{XX} = 0$, the defining property of a
cuscuton: zero propagating scalar degrees of freedom.  The GB correction
$\eps X$ breaks the exact cuscuton constraint, producing a sound speed
\begin{equation}
  c_s^2 = \frac{C_q}{\eps}\,,\qquad
  C_q \equiv \frac{2(1 - q_0)}{3|1 + q_0|}\,,
  \label{eq:cs}
\end{equation}
where $q_0$ is the deceleration parameter.  For $q_0 \in [-0.53, -0.39]$,
this gives $c_s \in [12c,\; 15c]$---superluminal but non-pathological.  The
Gauss--Bonnet correction reintroduces a propagating mode, so this is a genuine
front velocity rather than a phase-velocity artefact; causality survives because
the RS geometry supplies a preferred frame (the brane rest frame) in which the
Babichev--Mukhanov--Vikman closed-signal-curve construction cannot be carried
out~\cite{Babichev2008,Bruneton2007}.

The brane mass parameter $\mu^2$ is determined by the Goldberger--Wise (GW)
stabilisation mechanism~\cite{GW1999}.  The GW scalar with scaling dimension
$\Delta = 2 + \epsGW$ satisfies Robin boundary conditions whose eigenvalues
determine $\mu^2$ via the spectral action.  All ingredients except $\epsGW$
are derived from the spectral triple.  The algebraic chain is
\begin{equation}
  \mathbb{O} \xrightarrow{M_{\mathrm{oct}}}
  Y_f \xrightarrow{b_{3/2}}
  \alpha_{\mathrm{UV}} \xrightarrow{\text{Robin}}
  \mu^2(\epsGW) \xrightarrow{\text{JC}}
  w_0\,,
  \label{eq:algebraic-chain}
\end{equation}
where $M_{\mathrm{oct}}$ is the democratic mass matrix from the octonionic
mixing (eigenvalues $\{1/2, 1/2, 2\}$, $S_3$-symmetric), $b_{3/2}$ is the
fermion profile at the one-loop unification scale, $\alpha_{\mathrm{UV}}$ is
the UV coupling from the spectral action, and the junction conditions (JC) at
the IR brane produce the equation of state.

Fitting $\epsGW$ to the DESI~2024 BAO data~\cite{DESI2024} yields
$\epsGW = 0.275^{+0.072}_{-0.106}$ at $1\sigma$, giving:
\begin{equation}
  \boxed{w_0 = -0.830\,,\qquad w_a = 0 \;\text{(exactly)}\,.}
  \label{eq:w0-prediction}
\end{equation}
The $w_a = 0$ prediction is structural: the RS geometry produces a
constant equation of state at late times (the warp factor is static after
stabilisation).  This is testable by DESI~Year~5 at $3.8\sigma$ against the
$w_0$--$w_a$ posterior.


%% ============================================================
%%  THRESHOLD AND NUMERICAL VALIDATION
%% ============================================================

\textit{Topological threshold.}---The one-loop gauge coupling
threshold correction from the internal $\dP$ geometry involves the analytic
torsion of the K\"ahler--Einstein (KE) metric.  For a del~Pezzo surface of
degree~$d$ with canonical class self-intersection $K^2 = d$, the threshold
ratio is
\begin{equation}
  \frac{a_1}{a_2} = \frac{\ln K^2}{\sqrt{K^2 - 1}}
  = \frac{\ln 3}{\sqrt{2}} \approx 0.777\,,
  \label{eq:threshold}
\end{equation}
which depends only on the topological invariant $K^2 = 3$.  The existence of
the KE metric is guaranteed by the Tian--Yau theorem~\cite{TianYau1987}.

\textit{Numerical validation.}---We verify the KE geometry
using the Donaldson balanced metric algorithm~\cite{Donaldson2005} at embedding
degrees $k = 8, 12, 15$ on the Fermat cubic
$a^3 + b^3 + c^3 + 1 = 0 \subset \mathbb{CP}^2$.  At $k = 8$ (109 sections,
$(\mathbb{Z}/3)^2$ block decomposition, $\delta G = 5.3 \times 10^{-9}$), the
Galerkin eigenvalue computation with 196 basis functions on 600{,}000 sample
points yields $\lambda_1 = 1.461$, recovering 97.4\% of the KE Lichnerowicz
bound $\lambda_1 \geq 3/2$.

Eigenvector decomposition under $(\mathbb{Z}/3)^2$ reveals that $\lambda_1$
and $\lambda_2 = 1.494$ occupy different $S_3$ representations (standard
vs.\ trivial orbit), confirming their near-degeneracy is accidental.  The
$S_3$ splitting $\Delta\lambda_{S_3} = |\lambda_{(1,0)} - \lambda_{(0,1)}|$
shrinks by 31\% from $k = 8$ to $k = 15$ (Table~\ref{tab:splitting}),
confirming convergence toward the $S_3$-symmetric KE geometry.

\begin{table}[b]
\caption{$S_3$ splitting convergence across balanced-metric embedding
degrees.  $\lambda_{(q_1,q_2)}$ denotes the first eigenvalue in the
$(q_1, q_2)$ charge sector.}
\label{tab:splitting}
\begin{ruledtabular}
\begin{tabular}{ccccc}
$k$ & $N_{\text{sec}}$ & $\lambda_{(1,0)}$ & $\lambda_{(0,1)}$
& $\Delta\lambda_{S_3}$ \\
\hline
8  & 109 & 1.461 & 1.741 & 0.279 \\
12 & 235 & 1.482 & 1.695 & 0.214 \\
15 & 361 & 1.483 & 1.675 & 0.193 \\
\end{tabular}
\end{ruledtabular}
\end{table}


%% ============================================================
%%  DISCUSSION
%% ============================================================

\textit{Discussion.}---The framework presented here provides a
concrete UV completion of the quadratic gravity program.  Where
QQG~\cite{Afshordi2026} introduces $R^2$ corrections with adjustable
coefficients, the spectral action on the 5D warped orbifold \emph{derives}
these corrections with $\CGB = 2/3$, fixed by the Davis junction conditions.
The entire bosonic action---gravity, Standard~Model gauge fields, Higgs
sector, and dark energy---emerges from a single spectral triple with one
undetermined cosmological parameter ($\epsGW$).

The falsifiable predictions are: (i)~$w_0 = -0.830$, testable by DESI~Y5;
(ii)~$w_a = 0$ exactly, testable at $3.8\sigma$; (iii)~$c_s \in [12c, 15c]$,
constrainable by CMB lensing and large-scale structure;
(iv)~$\alpha_T = 0$ for gravitational-wave speed (GW170817-compatible).
The internal geometry prediction---$K^2(\dP) = 3$ selecting the threshold
via~\eqref{eq:threshold}---is topological and therefore exact.

Full derivations, the 170-page monograph including the fermion sector
(octonionic Yukawa couplings, neutrino masses, leptogenesis), and all
numerical computations are available in the companion
paper~\cite{IgguldenSchnell2026}.

\begin{acknowledgments}
The author thanks Clawd (AI collaborator) for sustained computational
assistance throughout this work.
\end{acknowledgments}


%% ============================================================
%%  REFERENCES
%% ============================================================

\begin{thebibliography}{99}

\bibitem{Afshordi2026}
R.~Liu, J.~Quintin, and N.~Afshordi,
Phys.\ Rev.\ Lett.\ \textbf{136}, 111501 (2026).

\bibitem{CC1997}
A.~H.~Chamseddine and A.~Connes,
Commun.\ Math.\ Phys.\ \textbf{186}, 731 (1997).

\bibitem{CC2010}
A.~H.~Chamseddine and A.~Connes,
J.\ Geom.\ Phys.\ \textbf{61}, 317 (2011).

\bibitem{CCM2007}
A.~H.~Chamseddine, A.~Connes, and M.~Marcolli,
Adv.\ Theor.\ Math.\ Phys.\ \textbf{11}, 991 (2007).

\bibitem{CC2012}
A.~H.~Chamseddine and A.~Connes,
J.\ High Energy Phys.\ \textbf{2012}(09), 104 (2012).

\bibitem{RS1999}
L.~Randall and R.~Sundrum,
Phys.\ Rev.\ Lett.\ \textbf{83}, 3370 (1999).

\bibitem{Davis2003}
S.~C.~Davis,
Phys.\ Rev.\ D \textbf{67}, 024030 (2003).

\bibitem{GW1999}
W.~D.~Goldberger and M.~B.~Wise,
Phys.\ Rev.\ Lett.\ \textbf{83}, 4922 (1999).

\bibitem{Afshordi2007}
N.~Afshordi, D.~J.~H.~Chung, and G.~Geshnizjani,
Phys.\ Rev.\ D \textbf{75}, 083513 (2007).

\bibitem{Babichev2008}
E.~Babichev, V.~Mukhanov, and A.~Vikman,
J.\ High Energy Phys.\ \textbf{2008}(02), 101 (2008).

\bibitem{Bruneton2007}
J.-P.~Bruneton,
Phys.\ Rev.\ D \textbf{75}, 085013 (2007).

\bibitem{DESI2024}
DESI Collaboration,
arXiv:2404.03002 [astro-ph.CO] (2024).

\bibitem{TianYau1987}
G.~Tian and S.-T.~Yau,
Commun.\ Math.\ Phys.\ \textbf{112}, 175 (1987).

\bibitem{Donaldson2005}
S.~K.~Donaldson,
J.\ Differ.\ Geom.\ \textbf{59}, 479 (2001).

\bibitem{Freudenthal1954}
H.~Freudenthal,
Indag.\ Math.\ \textbf{16}, 218 (1954).

\bibitem{Tits1966}
J.~Tits,
Indag.\ Math.\ \textbf{28}, 223 (1966).

\bibitem{IgguldenSchnell2026}
C.~W.~Iggulden-Schnell,
``Noncommutative spectral geometry on warped orbifolds:
topological couplings, gravitational corrections, and cosmological
predictions,'' Zenodo (2026).
%% Fill Zenodo DOI after upload

\end{thebibliography}

\end{document}
